% ৪.২ ওয়েবপেজ, ওয়েবসাইট ও হোমপেজ
\section{ওয়েবপেজ, ওয়েবসাইট ও হোমপেজ}

\begin{learningbox}
এই টপিক শেষে শিক্ষার্থী পারবে:
\begin{itemize}
\item ওয়েবপেজ, ওয়েবসাইট ও হোমপেজের সংজ্ঞা লিখতে।
\item ওয়েবপেজ ও ওয়েবসাইটের পার্থক্য ব্যাখ্যা করতে।
\item client ও server-এর ভূমিকা বুঝতে।
\item static ও dynamic webpage-এর মৌলিক ধারণা লিখতে।
\item শিক্ষা প্রতিষ্ঠান বা ব্যবসার জন্য website ব্যবহারের গুরুত্ব লিখতে।
\end{itemize}
\end{learningbox}

\subsection{ভূমিকা}

ওয়েব ডিজাইন শেখার আগে ওয়েবপেজ, ওয়েবসাইট ও হোমপেজ - এই তিনটি ধারণা পরিষ্কারভাবে বোঝা জরুরি। কারণ HTML ব্যবহার করে মূলত ওয়েবপেজ তৈরি করা হয়, আর একাধিক ওয়েবপেজ মিলেই একটি পূর্ণাঙ্গ ওয়েবসাইট গঠিত হয়।

একটি বই যেমন অনেকগুলো পৃষ্ঠা নিয়ে তৈরি হয়, তেমনি একটি ওয়েবসাইটও অনেকগুলো ওয়েবপেজ নিয়ে তৈরি হয়। বইয়ের প্রথম পৃষ্ঠা বা সূচিপত্র যেমন পাঠককে পুরো বইয়ের ধারণা দেয়, তেমনি ওয়েবসাইটের হোমপেজ user-কে পুরো website সম্পর্কে প্রথম ধারণা দেয়।

\subsection{ওয়েবপেজ কী?}

\begin{definitionbox}{ওয়েবপেজ}
ইন্টারনেট ব্যবহারকারীদের দেখার জন্য ওয়েব সার্ভারে সংরক্ষিত HTML বা অন্যান্য ওয়েব প্রযুক্তি দ্বারা তৈরি স্বতন্ত্র ওয়েব ডকুমেন্টকে ওয়েবপেজ বলে।
\end{definitionbox}

সহজভাবে বলা যায়, browser-এ যে একক page দেখা যায় সেটিই ওয়েবপেজ। একটি ওয়েবপেজে text, image, hyperlink, list, table, audio, video, form ইত্যাদি থাকতে পারে।

\begin{examplebox}{উদাহরণ}
একটি কলেজের ওয়েবসাইটে Home page, About page, Notice page, Admission page, Result page এবং Contact page থাকতে পারে। এগুলোর প্রতিটি আলাদা আলাদা ওয়েবপেজ।
\end{examplebox}

\begin{keypointbox}{ওয়েবপেজের বৈশিষ্ট্য}
\begin{itemize}
\item এটি browser-এর মাধ্যমে দেখা যায়।
\item সাধারণত HTML দিয়ে তৈরি হয়।
\item এতে text, image, link, table, list, audio, video বা form থাকতে পারে।
\item প্রতিটি ওয়েবপেজের আলাদা URL থাকতে পারে।
\item এটি web server-এ সংরক্ষিত থাকতে পারে।
\item user internet ব্যবহার করে ওয়েবপেজ দেখতে পারে।
\item একটি ওয়েবপেজ একই website-এর অন্য page-এর সাথে hyperlink দ্বারা যুক্ত থাকতে পারে।
\end{itemize}
\end{keypointbox}

\subsection{ওয়েবসাইট কী?}

\begin{definitionbox}{ওয়েবসাইট}
একই domain name-এর অধীনে সংরক্ষিত ও পরস্পর সম্পর্কযুক্ত একাধিক ওয়েবপেজের সমষ্টিকে ওয়েবসাইট বলে।
\end{definitionbox}

অর্থাৎ, একটি website হলো অনেকগুলো linked webpage-এর সংগঠিত collection। Website সাধারণত কোনো ব্যক্তি, প্রতিষ্ঠান, ব্যবসা, শিক্ষা প্রতিষ্ঠান, সংবাদমাধ্যম বা সরকারি সেবার তথ্য উপস্থাপন করে।

\begin{examplebox}{উদাহরণ}
\texttt{www.educationboard.gov.bd} একটি ওয়েবসাইট। এর মধ্যে Result page, Notice page, Circular page, Contact page এবং Login page থাকতে পারে। এগুলো একসাথে মিলে ওয়েবসাইট গঠন করে।
\end{examplebox}

\subsection{হোমপেজ কী?}

\begin{definitionbox}{হোমপেজ}
কোনো ওয়েবসাইটে প্রবেশ করার পর প্রথম যে ওয়েবপেজটি প্রদর্শিত হয় তাকে হোমপেজ বলে।
\end{definitionbox}

হোমপেজকে website-এর প্রধান বা starting page বলা যায়। সাধারণত এখান থেকেই user website-এর অন্যান্য page-এ যায়। অনেক website-এ home page-এর file name \texttt{index.html} হয়।

\begin{keypointbox}{হোমপেজে সাধারণত যা থাকে}
\begin{itemize}
\item ওয়েবসাইটের নাম বা logo।
\item প্রধান menu বা navigation bar।
\item গুরুত্বপূর্ণ notice বা update।
\item ওয়েবসাইটের সংক্ষিপ্ত পরিচিতি।
\item অন্যান্য page-এর link।
\item search box বা contact information।
\item image, banner বা featured content।
\end{itemize}
\end{keypointbox}

\subsection{ওয়েবপেজ, ওয়েবসাইট ও হোমপেজের সম্পর্ক}

সম্পর্কটি সহজভাবে মনে রাখা যায়:

\begin{center}
\begin{tikzpicture}[
  root/.style={draw, rounded corners, fill=green!10, minimum width=3.15cm, minimum height=0.82cm, align=center, font=\small\bfseries},
  page/.style={draw, rounded corners, fill=blue!6, minimum width=2.35cm, minimum height=0.72cm, align=center, font=\scriptsize},
  note/.style={font=\scriptsize, align=center},
  arrow/.style={->, thick}
]
\node[root] (site) {ABC College Website};
\node[page, below left=1.05cm and 2.2cm of site] (home) {Home page\\\texttt{index.html}};
\node[page, below=1.05cm of site] (notice) {Notice page\\\texttt{notice.html}};
\node[page, below right=1.05cm and 2.2cm of site] (contact) {Contact page\\\texttt{contact.html}};
\node[page, below=0.8cm of home] (about) {About page};
\node[page, below=0.8cm of contact] (result) {Result page};
\draw[arrow] (site) -- (home);
\draw[arrow] (site) -- (notice);
\draw[arrow] (site) -- (contact);
\draw[arrow] (home) -- (about);
\draw[arrow] (contact) -- (result);
\node[note, above=0.1cm of home] {প্রথম/প্রধান page};
\node[note, below=0.1cm of notice] {প্রতিটি box একটি webpage};
\end{tikzpicture}\\[3pt]
\textit{চিত্র: অনেকগুলো webpage মিলে website, আর প্রথম page হলো homepage}
\end{center}

এখানে Home page-ও একটি ওয়েবপেজ, কিন্তু এটি ওয়েবসাইটের প্রথম ও প্রধান পেজ। অন্য page গুলো website-এর information section হিসেবে কাজ করে।

\subsection{ওয়েবপেজ ও ওয়েবসাইটের পার্থক্য}

\begin{center}
\begin{tabular}{|p{0.22\textwidth}|p{0.31\textwidth}|p{0.31\textwidth}|}
\hline
\textbf{বিষয়} & \textbf{ওয়েবপেজ} & \textbf{ওয়েবসাইট} \\
\hline
অর্থ & একক web document & একাধিক webpage-এর সমষ্টি \\
\hline
গঠন & একটি page & অনেকগুলো page \\
\hline
URL & একটি নির্দিষ্ট URL থাকতে পারে & একটি domain-এর অধীনে অনেক URL থাকে \\
\hline
উদাহরণ & Contact page, Notice page & কলেজের পুরো website \\
\hline
সম্পর্ক & ওয়েবসাইটের অংশ & ওয়েবপেজগুলোকে একত্র করে তৈরি \\
\hline
ব্যবহার & নির্দিষ্ট তথ্য দেখায় & সম্পূর্ণ প্রতিষ্ঠান, ব্যক্তি বা service উপস্থাপন করে \\
\hline
\end{tabular}
\end{center}

\subsection{ওয়েবসাইট তৈরি করতে কী কী লাগে?}

একটি ওয়েবসাইট তৈরি ও প্রকাশ করতে সাধারণত কয়েকটি ধাপ অনুসরণ করা হয়।

\begin{enumerate}
\item \textbf{Domain name registration:} website-এর address নেওয়া হয়, যেমন \texttt{example.com} বা \texttt{college.edu.bd}।
\item \textbf{Webpage design:} HTML, CSS এবং প্রয়োজন অনুযায়ী image, table, form, link ব্যবহার করে page তৈরি করা হয়।
\item \textbf{Hosting বা server:} website-এর HTML, image, CSS, JavaScript file server-এ রাখা হয়।
\item \textbf{DNS configuration:} domain name যেন সঠিক server-এর সাথে যুক্ত হয় তা নির্ধারণ করা হয়।
\item \textbf{Testing:} browser দিয়ে link, image, table, form ও page layout ঠিক আছে কি না দেখা হয়।
\item \textbf{Search engine submission/SEO:} user যেন search engine-এ website খুঁজে পায়, সে জন্য basic optimization করা হয়।
\end{enumerate}

\begin{boardbox}{পরীক্ষার জন্য}
ওয়েবসাইট তৈরি করতে domain name, webpage design, hosting/server, DNS configuration এবং testing দরকার হয়। HSC উত্তরে সংক্ষেপে domain, hosting ও webpage design লিখলেই মূল ধারণা পরিষ্কার হয়।
\end{boardbox}

\subsection{ওয়েবসাইটের প্রধান দুই অংশ: Client ও Server}

ওয়েবসাইট ব্যবহারের সময় প্রধানত দুইটি অংশ কাজ করে - client এবং server।

\subsubsection{Client}

ব্যবহারকারীর device বা browser-কে client বলা হয়। যেমন laptop, mobile, tablet এবং Chrome, Firefox বা Edge browser। Client server-এর কাছে webpage চেয়ে request পাঠায়।

\subsubsection{Server}

যে computer বা system-এ website-এর file সংরক্ষিত থাকে এবং client-এর request অনুযায়ী webpage পাঠায় তাকে server বলা হয়। Server HTML, CSS, image, video, script বা database response পাঠাতে পারে।

\begin{center}
\begin{tikzpicture}[
  box/.style={draw, rounded corners, fill=green!6, minimum width=3cm, minimum height=0.82cm, align=center, font=\small},
  arrow/.style={->, thick}
]
\node[box] (client) {Client\\Browser/Device};
\node[box, right=2.3cm of client] (server) {Server\\Website files};
\draw[arrow] ([yshift=0.16cm]client.east) -- node[above, font=\scriptsize]{request} ([yshift=0.16cm]server.west);
\draw[arrow] ([yshift=-0.16cm]server.west) -- node[below, font=\scriptsize]{response} ([yshift=-0.16cm]client.east);
\node[font=\scriptsize, align=center, below=0.45cm of client] {User এখানে page দেখে};
\node[font=\scriptsize, align=center, below=0.45cm of server] {HTML, image, CSS রাখা থাকে};
\end{tikzpicture}\\[3pt]
\textit{চিত্র: Client request পাঠায়, server response পাঠায়}
\end{center}

\begin{boardbox}{সৃজনশীল প্রশ্নে লেখার মতো}
Client হলো ব্যবহারকারীর browser বা device, যা server-এ request পাঠায়। Server হলো এমন computer, যেখানে ওয়েবসাইটের file সংরক্ষিত থাকে এবং client-এর request অনুযায়ী webpage পাঠায়। এই client-server communication-এর মাধ্যমেই user website দেখতে পারে।
\end{boardbox}

\subsection{ওয়েবপেজের প্রকারভেদ}

HSC level-এ ওয়েবপেজ সাধারণত দুই ধরনের হিসেবে আলোচনা করা হয় - static webpage এবং dynamic webpage।

\subsubsection{Static Webpage}

\begin{definitionbox}{Static Webpage}
যে ওয়েবপেজের content সাধারণত fixed থাকে এবং ব্যবহারকারীভেদে পরিবর্তিত হয় না তাকে static webpage বলে।
\end{definitionbox}

\begin{keypointbox}{Static webpage-এর বৈশিষ্ট্য}
\begin{itemize}
\item HTML/CSS দিয়ে সহজে তৈরি করা যায়।
\item server-side processing সাধারণত লাগে না।
\item content manually update করতে হয়।
\item দ্রুত load হয়।
\item ছোট website, portfolio, notice page বা পরিচিতি page-এর জন্য উপযোগী।
\end{itemize}
\end{keypointbox}

\begin{examplebox}{উদাহরণ}
কলেজের পরিচিতি page, সাধারণ notice page, personal portfolio page ইত্যাদি static webpage হতে পারে।
\end{examplebox}

\subsubsection{Dynamic Webpage}

\begin{definitionbox}{Dynamic Webpage}
যে ওয়েবপেজ user input, সময়, database বা program logic অনুযায়ী পরিবর্তিত হয় তাকে dynamic webpage বলে।
\end{definitionbox}

\begin{keypointbox}{Dynamic webpage-এর বৈশিষ্ট্য}
\begin{itemize}
\item server-side বা client-side script ব্যবহার হতে পারে।
\item database থেকে তথ্য দেখাতে পারে।
\item login, search, result, comment বা order system থাকতে পারে।
\item ব্যবহারকারীভেদে content পরিবর্তিত হতে পারে।
\item static webpage-এর তুলনায় design ও maintenance বেশি জটিল হতে পারে।
\end{itemize}
\end{keypointbox}

\begin{examplebox}{উদাহরণ}
Facebook news feed, online result search page, e-commerce product page, student login portal ইত্যাদি dynamic webpage-এর উদাহরণ।
\end{examplebox}

\subsection{ওয়েবসাইটের সুবিধা}

একটি website ব্যক্তি, প্রতিষ্ঠান বা organization-এর তথ্য দ্রুত ও সহজভাবে সবার কাছে পৌঁছে দিতে পারে।

\begin{itemize}
\item তথ্য দ্রুত প্রকাশ করা যায়।
\item ২৪ ঘণ্টা যেকোনো স্থান থেকে তথ্য দেখা যায়।
\item প্রতিষ্ঠান বা ব্যবসার প্রচার সহজ হয়।
\item কম খরচে বিশ্বব্যাপী তথ্য ছড়ানো যায়।
\item online service দেওয়া যায়।
\item user সহজে notice, result, admission information বা routine পেতে পারে।
\item image, video, document, form ও download link যুক্ত করা যায়।
\item customer, student বা visitor সহজে যোগাযোগ করতে পারে।
\end{itemize}

\subsection{ওয়েবসাইটের ব্যবহার}

\begin{center}
\begin{tabular}{|p{0.24\textwidth}|p{0.58\textwidth}|}
\hline
\textbf{ক্ষেত্র} & \textbf{ব্যবহার} \\
\hline
শিক্ষা & online class, notice, admission form, result, syllabus, routine \\
\hline
ব্যবসা & product display, online order, customer support, payment, brand promotion \\
\hline
সরকারি কাজ & নাগরিক সেবা, birth registration, passport service, tax information, circular \\
\hline
ব্যক্তিগত কাজ & portfolio, blog, CV website, project showcase \\
\hline
সংবাদমাধ্যম & online news, article, image, video, archive \\
\hline
\end{tabular}
\end{center}

\subsection{বইয়ে দেওয়ার মতো সুন্দর উদাহরণ}

ধরা যাক, \textbf{ABC College} নামে একটি কলেজের ওয়েবসাইট তৈরি করা হবে।

\begin{center}
\begin{tabular}{|p{0.30\textwidth}|p{0.45\textwidth}|}
\hline
\textbf{File/Page} & \textbf{কাজ} \\
\hline
\texttt{index.html} & Home page \\
\hline
\texttt{about.html} & কলেজ পরিচিতি \\
\hline
\texttt{notice.html} & Notice board \\
\hline
\texttt{result.html} & Result page \\
\hline
\texttt{gallery.html} & ছবি বা gallery \\
\hline
\texttt{contact.html} & যোগাযোগ \\
\hline
\end{tabular}
\end{center}

এখানে প্রতিটি \texttt{.html} file হলো একটি ওয়েবপেজ। সবগুলো page একই domain-এর অধীনে থাকলে সেটিই হবে ABC College-এর ওয়েবসাইট। \texttt{index.html} সাধারণত home page হিসেবে কাজ করে।

\subsection{পরীক্ষার জন্য সংক্ষিপ্ত উত্তর}

\subsubsection{প্রশ্ন ১: ওয়েবপেজ কী?}

\begin{boardbox}{উত্তর}
ইন্টারনেট ব্যবহারকারীদের দেখার জন্য ওয়েব সার্ভারে সংরক্ষিত HTML বা অন্যান্য ওয়েব প্রযুক্তি দ্বারা তৈরি স্বতন্ত্র ওয়েব ডকুমেন্টকে ওয়েবপেজ বলে। একটি ওয়েবসাইটের প্রতিটি আলাদা page-ই ওয়েবপেজ হিসেবে বিবেচিত হয়।
\end{boardbox}

\subsubsection{প্রশ্ন ২: ওয়েবসাইট কী?}

\begin{boardbox}{উত্তর}
একই domain name-এর অধীনে সংরক্ষিত ও পরস্পর সম্পর্কযুক্ত একাধিক ওয়েবপেজের সমষ্টিকে ওয়েবসাইট বলে। ওয়েবসাইটে সাধারণত home page, contact page, notice page, image page ইত্যাদি থাকে।
\end{boardbox}

\subsubsection{প্রশ্ন ৩: হোমপেজ কী?}

\begin{boardbox}{উত্তর}
কোনো ওয়েবসাইটে প্রবেশ করার পর প্রথম যে ওয়েবপেজটি প্রদর্শিত হয় তাকে হোমপেজ বলে। এটি ওয়েবসাইটের প্রধান page, যেখান থেকে অন্যান্য page-এ যাওয়ার link থাকে।
\end{boardbox}

\subsubsection{প্রশ্ন ৪: ওয়েবসাইটের দুটি অংশ ব্যাখ্যা কর।}

\begin{boardbox}{উত্তর}
ওয়েবসাইট ব্যবহারের ক্ষেত্রে প্রধান দুটি অংশ হলো client ও server। Client হলো ব্যবহারকারীর browser বা device, যা server-এ request পাঠায়। Server হলো এমন computer, যেখানে ওয়েবসাইটের file সংরক্ষিত থাকে এবং client-এর request অনুযায়ী webpage পাঠায়।
\end{boardbox}

\subsection{Board/Exam Focus}

\begin{keypointbox}{মনে রাখার point}
\begin{itemize}
\item ওয়েবপেজ হলো একক web document।
\item ওয়েবসাইট হলো একই domain-এর অধীনে একাধিক webpage-এর সমষ্টি।
\item হোমপেজ হলো website-এর প্রথম বা প্রধান page।
\item ওয়েবসাইটের file web server-এ থাকে।
\item browser client হিসেবে server থেকে webpage request করে।
\item HTML web page-এর structure তৈরি করার প্রধান language।
\item Static webpage fixed থাকে; dynamic webpage user/database/input অনুযায়ী পরিবর্তিত হতে পারে।
\end{itemize}
\end{keypointbox}

\subsection{Practice MCQ}

\begin{enumerate}
\item একই domain-এর অধীনে একাধিক webpage-এর সমষ্টিকে কী বলে?\\
ক) Web browser \quad খ) Website \quad গ) Search engine \quad ঘ) Web server\\
\textbf{উত্তর:} খ) Website

\item কোনো website-এ প্রবেশ করলে প্রথম যে page দেখা যায় তাকে কী বলে?\\
ক) Search page \quad খ) Server page \quad গ) Home page \quad ঘ) Link page\\
\textbf{উত্তর:} গ) Home page

\item ওয়েবসাইটের file কোথায় সংরক্ষিত থাকে?\\
ক) Keyboard \quad খ) Web server \quad গ) Printer \quad ঘ) Scanner\\
\textbf{উত্তর:} খ) Web server

\item ওয়েবপেজ তৈরির মৌলিক ভাষা কোনটি?\\
ক) HTML \quad খ) SQL \quad গ) C \quad ঘ) Python\\
\textbf{উত্তর:} ক) HTML

\item যে ওয়েবপেজ user input বা database অনুযায়ী পরিবর্তিত হয় তাকে কী বলে?\\
ক) Static webpage \quad খ) Dynamic webpage \quad গ) Printed page \quad ঘ) Offline page\\
\textbf{উত্তর:} খ) Dynamic webpage
\end{enumerate}

\subsection{CQ Practice}

\begin{practicebox}
\textbf{উদ্দীপক:} রাফি তার কলেজের জন্য একটি ওয়েবসাইট তৈরি করল। ওয়েবসাইটটিতে Home, About, Notice, Result এবং Contact নামে পাঁচটি page আছে। ব্যবহারকারী ওয়েবসাইটে প্রবেশ করলে প্রথমে Home page দেখা যায়। Notice page-এ কলেজের নতুন তথ্য প্রকাশ করা হয়।

\begin{enumerate}
\item ক. ওয়েবপেজ কী?
\item খ. হোমপেজকে ওয়েবসাইটের প্রধান page বলা হয় কেন?
\item গ. উদ্দীপকে উল্লিখিত পাঁচটি page-এর ভিত্তিতে ওয়েবসাইটের গঠন ব্যাখ্যা কর।
\item ঘ. শিক্ষা প্রতিষ্ঠানের জন্য ওয়েবসাইট ব্যবহারের গুরুত্ব বিশ্লেষণ কর।
\end{enumerate}
\end{practicebox}

\begin{boardbox}{সমাধান নির্দেশনা}
\textbf{ক.} ওয়েব সার্ভারে সংরক্ষিত browser-এ প্রদর্শনযোগ্য স্বতন্ত্র web document হলো ওয়েবপেজ।\\
\textbf{খ.} হোমপেজে সাধারণত website-এর পরিচিতি, menu ও অন্যান্য page-এর link থাকে, তাই এটি প্রধান page।\\
\textbf{গ.} Home, About, Notice, Result, Contact - প্রতিটি আলাদা ওয়েবপেজ; সবগুলো একত্রে একই domain-এর অধীনে থাকলে সেটি একটি ওয়েবসাইট।\\
\textbf{ঘ.} শিক্ষা প্রতিষ্ঠানে notice, result, admission info, routine, syllabus ও contact information দ্রুত প্রকাশের জন্য ওয়েবসাইট গুরুত্বপূর্ণ।
\end{boardbox}

\subsection{Common Mistake}

\begin{itemize}
\item ওয়েবপেজ ও ওয়েবসাইটকে একই মনে করা।
\item হোমপেজকে আলাদা website মনে করা।
\item domain name ও website address গুলিয়ে ফেলা।
\item server ও browser-এর কাজ এক মনে করা।
\item static ও dynamic webpage-এর পার্থক্য না বোঝা।
\item \texttt{index.html} page-কে শুধু file মনে করা, home page হিসেবে তার ভূমিকা না বোঝা।
\end{itemize}

\subsection{মনে রাখবে}

\begin{recapbox}
\begin{itemize}
\item ওয়েবপেজ হলো একটি একক web document।
\item ওয়েবসাইট হলো একাধিক linked webpage-এর সমষ্টি।
\item হোমপেজ হলো ওয়েবসাইটের প্রথম বা প্রধান page।
\item ওয়েবসাইট server-এ সংরক্ষিত থাকে।
\item browser server থেকে page এনে ব্যবহারকারীর সামনে দেখায়।
\item HTML ওয়েবপেজের structure তৈরি করার প্রধান language।
\end{itemize}
\end{recapbox}
